% SOURCE AUTHORITY:
% expose_VII_body.tex lines 421--460; scan indices 165--167; source folios 154--156.
% Hard stop before source line 465, section 2.

1.4. \underline{Extensions of a free commutative group}

Let $I$ be an object of $\underline T$, and put
\begin{equation}\tag{1.4.1}
 P=\mathbb Z[I],
\end{equation}
the free $\mathbb Z$-module generated by $I$ (SGA 4, IV, 11). With
every extension $E$ of $P$ by a given $\mathbb Z$-module $G$ (from now
on only commutative group extensions will occur), associate the
corresponding torsor under $G_P$, and then restrict it to $I$ along the
canonical morphism
\[
 I\longrightarrow P=\mathbb Z[I].
\]

%%%%%%%%%%  scan idx 166  |  folio 155  %%%%%%%%%%

This gives a canonical functor
\begin{equation}\tag{1.4.2}
 \operatorname{EXT}(P,G)\longrightarrow
 \operatorname{TORS}(I,G_I).
\end{equation}
I claim that this functor is an \underline{equivalence of categories}.
Equivalently, for every $G$, the maps
\begin{equation}\tag{1.4.3}
 \operatorname{Hom}_{\mathbb Z}(P,G)\longrightarrow
 H^0(I,G_I)=G(I)
\end{equation}
and
\begin{equation}\tag{1.4.4}
 \operatorname{Ext}^1_{\mathbb Z}(P,G)\longrightarrow H^1(I,G_I)
\end{equation}
are isomorphisms. More generally, if $A$ is a ring (not necessarily
commutative) of $\underline T$ and $P=A[I]$ is the free $A$-module
generated by $I$, there are canonical isomorphisms, functorial in the
variable $A$-module $G$,
\begin{equation}\tag{1.4.5}
 \operatorname{Ext}^i_A(A[I],G)\xrightarrow{\ \sim\ }H^i(I,G_I).
\end{equation}
For $A=\mathbb Z$ and $i=0,1$, these coincide with (1.4.3) and
(1.4.4). To define (1.4.5), recall that its left-hand side is the
derived functor of
\[
 G\longmapsto\operatorname{Hom}_A(A[I],G),
\]
which, by the definition of $A[I]$, is isomorphic to
$G\longmapsto G(I)=H^0(I,G_I)$. Since $G\longmapsto G_I$ carries
injective modules to injective modules of $\underline T_{/I}$ (SGA 4,
V), the assertion follows. We leave to the reader the verification of
the stated compatibility with (1.4.3) and (1.4.4).

Finally, the functor (1.4.2) is compatible with inverse images under a
morphism of topoi
\[
 f:\underline T'\longrightarrow\underline T.
\]

